% Spettro elettromagnetico: regioni in ordine di frequenza crescente
\begin{tikzpicture}[>=Stealth, font=\small,
  regione/.style={minimum width=2.1cm, minimum height=1cm, align=center,
                  text width=1.9cm, inner sep=1pt}]
  \foreach \i/\nome/\col in {
      0/{onde\\radio}/formula!20,
      1/microonde/formula!35,
      2/infrarosso/attenzione!30,
      4/ultravioletto/approfondimento!25,
      5/{raggi X}/approfondimento!45,
      6/{raggi $\gamma$}/approfondimento!65} {
    \node[regione, fill=\col] at (2.1*\i,0) {\nome};
  }
  % il visibile, sfumato dal rosso al violetto
  \shade[left color=red!60, right color=violet!60, middle color=green!50]
    (2.1*3-1.05,-0.5) rectangle (2.1*3+1.05,0.5);
  \node[regione] at (2.1*3,0) {visibile};
  % frecce
  \draw[->, thick] (-1.05,0.85) -- (13.65,0.85)
    node[pos=0.28, above] {frequenza ed energia del fotone crescenti};
  \draw[<-, thick] (-1.05,-0.85) -- (13.65,-0.85)
    node[pos=0.5, below] {lunghezza d'onda crescente};
  % soglia di ionizzazione, nell'ultravioletto
  \draw[attenzione, very thick, dashed] (2.1*4+0.5,-0.6) -- (2.1*4+0.5,1.3);
  \node[attenzione, above] at (2.1*4+0.5,1.3)
    {soglia di ionizzazione (circa \SI{10}{\electronvolt})};
  \node at (2.1*1.5,-1.75) {$\longleftarrow$ non ionizzanti};
  \node at (2.1*5.6,-1.75) {ionizzanti $\longrightarrow$};
\end{tikzpicture}
